\begin{table}[!htbp]
\setlength{\abovecaptionskip}{0pt}
\centering
\caption{Heterogeneity in Procurement Effects by Plaintiff Type}
\label{tab:admin_plaintiff_heterogeneity}
\begin{threeparttable}
\begin{tabular}{lcccc}
\toprule
 & (1) & (2) & (3) & (4) \\
Subsample & Entity plaintiff & Entity plaintiff & Individual plaintiff & Individual plaintiff \\
\midrule
Treatment $\times$ Post & 0.037$^{***}$ & 0.040$^{***}$ & 0.030$^{***}$ & 0.034$^{***}$ \\
& (0.006) & (0.006) & (0.005) & (0.005) \\
\addlinespace
Opposing counsel &  & -0.030$^{***}$ &  & -0.038$^{***}$ \\
&  & (0.002) &  & (0.002) \\
\addlinespace
Observations & 275,406 & 275,406 & 383,765 & 383,765 \\
$R^2$ & 0.043 & 0.044 & 0.044 & 0.046 \\
Equality-test p-value & \multicolumn{2}{c}{Cols. 1 = 3: <0.001} & \multicolumn{2}{c}{Cols. 2 = 4: <0.001} \\
Opposing counsel &  & Yes &  & Yes \\
Court Fixed Effects & Yes & Yes & Yes & Yes \\
Year Fixed Effects & Yes & Yes & Yes & Yes \\
Cause-Group Fixed Effects & Yes & Yes & Yes & Yes \\
\bottomrule
\end{tabular}
\begin{tablenotes}[flushleft]
\footnotesize
\item \textit{Note:} Linear-probability coefficients with the government-win indicator as the outcome, estimated on the indicated plaintiff sub-sample. Entity-plaintiff sub-sample $N=275,406$; individual-plaintiff sub-sample $N=383,765$. Even columns add an indicator for whether the opposing party appears with counsel. The equality-test row reports two-sided $p$-values for equality of the procurement coefficient between entity- and individual-plaintiff cases, estimated from the corresponding pooled interaction specification. Standard errors clustered by city and court. $^{*}p<0.10$, $^{**}p<0.05$, $^{***}p<0.01$.
\end{tablenotes}
\end{threeparttable}
\end{table}
